% =============================================================================
% Partie 3 — Architecture de la base de données et DDL
% =============================================================================
\chapter{Architecture de la base et DDL}

\epigraphe{On ne construit pas une maison en commençant par le toit, ni une base en commençant par les requêtes.}

\begin{imagine}
Un architecte ne commence jamais par poser des briques. Il dessine d'abord le plan : les pièces, les
portes qui les relient, ce qui porte le poids de la maison. Une erreur sur le plan coûte un coup de
gomme ; la même erreur une fois le mur construit coûte un chantier.

Pour une base de données, c'est pareil. Le \textbf{modèle conceptuel} est le croquis, le
\textbf{modèle relationnel} est le plan coté, et le \textbf{code DDL} est la construction.
\end{imagine}

% -----------------------------------------------------------------------------
\section{Du besoin au modèle conceptuel}

On part des phrases du métier et on repère les \textbf{entités} (les « choses » dont on garde la trace)
et les \textbf{associations} entre elles :
\begin{itemize}
\item un \textbf{client} demande un ou plusieurs \textbf{prêts} ; chaque prêt appartient à un seul client ;
\item un prêt accordé comporte plusieurs \textbf{échéances} ; un prêt refusé n'en a aucune ;
\item un prêt est évalué par un ou plusieurs \textbf{modèles}, chaque évaluation est une \textbf{prédiction} ;
\item un prêt appartient à l'échantillon d'apprentissage ou à l'échantillon de test du modèle ;
\item chaque changement de statut d'un prêt est \textbf{journalisé}, et chaque fin de mois produit une
\textbf{provision} par prêt.
\end{itemize}

La figure~\ref{fig:mcd} traduit ces phrases en \textbf{modèle conceptuel de données} (MCD, notation
Merise). Les \textbf{cardinalités} se lisent ainsi : « un client demande au minimum 1 et au maximum
$n$ prêts ; un prêt est demandé par exactement 1 client ».

\begin{figure}[ht]
\centering
\begin{tikzpicture}[font=\small,
  entite/.style={rectangle split, rectangle split parts=2, draw=bleuML, line width=0.7pt,
     rectangle split part fill={bleuFond, white}, text width=2.95cm, align=center, inner sep=3pt},
  asso/.style={ellipse, draw=vertCadre, fill=vertFond, line width=0.6pt, minimum width=2cm, minimum height=0.8cm, align=center},
  card/.style={font=\footnotesize\bfseries, fill=white, inner sep=1pt}]
  \node[entite] (cli) at (0,0) {\textbf{CLIENT}\nodepart{two}\footnotesize âge, revenu,\\ logement, ancienneté};
  \node[asso] (dem) at (3.25,0) {demander};
  \node[entite] (pret) at (6.5,0) {\textbf{PRÊT}\nodepart{two}\footnotesize montant, durée,\\ taux, motif,\\ note, statuts};
  \node[asso] (comp) at (9.75,0) {comporter};
  \node[entite] (ech) at (12.8,0) {\textbf{ÉCHÉANCE}\nodepart{two}\footnotesize date, montant dû,\\ capital, payé};
  \node[asso] (eval) at (6.5,-2.6) {évaluer};
  \node[entite] (pred) at (6.5,-5) {\textbf{PRÉDICTION}\nodepart{two}\footnotesize version, PD, seuil,\\ décision, SHAP};
  \node[asso] (aff) at (3.0,-2.6) {affecter};
  \node[entite] (ech2) at (0,-5) {\textbf{ÉCHANTILLON}\nodepart{two}\footnotesize apprentissage\\ ou test};
  \node[asso] (prov) at (10.1,-2.6) {provisionner};
  \node[entite] (pm) at (12.8,-5) {\textbf{PROVISION}\nodepart{two}\footnotesize arrêté, stage,\\ EAD, PD, ECL};
  \draw (cli) -- node[card, pos=0.3] {1,n} (dem) -- node[card, pos=0.7] {1,1} (pret);
  \draw (pret) -- node[card, pos=0.3] {0,n} (comp) -- node[card, pos=0.7] {1,1} (ech);
  \draw (pret) -- node[card, pos=0.4] {0,n} (eval) -- node[card, pos=0.6] {1,1} (pred);
  \draw (pret) -- node[card, pos=0.35] {0,1} (aff) -- node[card, pos=0.6] {0,n} (ech2);
  \draw (pret) -- node[card, pos=0.35] {0,n} (prov) -- node[card, pos=0.6] {1,1} (pm);
\end{tikzpicture}
\caption{Modèle conceptuel de données (simplifié). Le journal des statuts suit le même schéma que les
provisions : un prêt en a de 0 à $n$.}
\label{fig:mcd}
\end{figure}

\subsection{Le passage au modèle relationnel}

Une seule règle suffit pour passer du croquis au plan :

\begin{regle}[Règle de passage]
Une association « un à plusieurs » devient une \textbf{clé étrangère} placée \textbf{du côté
« plusieurs »}. Un prêt a plusieurs échéances : c'est donc l'échéance qui porte la colonne
\code{pret_id}, et non le prêt qui porterait une liste d'échéances.
\end{regle}

On obtient le \textbf{modèle relationnel} de la figure~\ref{fig:mld} : chaque entité devient une table,
chaque association devient une clé étrangère. On y ajoute une table technique,
\code{parametres_risque}, qui centralise les paramètres métier.

\begin{figure}[p]
\centering
\begin{tikzpicture}[font=\footnotesize, >=Stealth,
  table/.style={rectangle split, rectangle split parts=2, draw=bleuML, line width=0.7pt,
     rectangle split part fill={bleuML, white}, text width=4.45cm, align=left, inner sep=3.5pt},
  tabled/.style={table, draw=grisTitre, dashed, rectangle split part fill={grisTitre, white}},
  card/.style={font=\scriptsize\bfseries, fill=white, inner sep=1pt}]
  \node[table, anchor=north] (cli) at (0,0) {\textcolor{white}{\textbf{clients}}\nodepart{two}
    \textbf{client\_id} \hfill PK\\ nom\_complet\\ email \hfill UNIQUE\\ age\\ revenu\_annuel\\
    statut\_logement\\ anciennete\_emploi\\ anciennete\_credit\\ defaut\_anterieur\\ date\_inscription};
  \node[table, anchor=north] (pret) at (5.5,0) {\textcolor{white}{\textbf{demandes\_pret}}\nodepart{two}
    \textbf{pret\_id} \hfill PK\\ \textit{client\_id} \hfill FK\\ montant\_demande\\ duree\_mois\\
    taux\_interet, taux\_impute\\ motif\_pret\\ note\_risque\\ date\_demande\\ date\_decision\\
    date\_premiere\_echeance\\ statut\_dossier\\ statut\_pret};
  \node[table, anchor=north] (ech) at (11.0,0) {\textcolor{white}{\scriptsize\textbf{echeancier\_remboursements}}\nodepart{two}
    \textbf{echeance\_id} \hfill PK\\ \textit{pret\_id} \hfill FK\\ num\_echeance\\ date\_echeance\\
    montant\_du\\ part\_capital\\ part\_interets\\ capital\_restant\_du\\ montant\_paye\\ date\_paiement\\[2pt]
    {\scriptsize UNIQUE (pret\_id, num\_echeance)}};
  \node[table, anchor=north] (pred) at (5.5,-7.3) {\textcolor{white}{\textbf{predictions\_scoring\_ml}}\nodepart{two}
    \textbf{prediction\_id} \hfill PK\\ \textit{pret\_id} \hfill FK\\ version\_modele\\ est\_champion\\
    probabilite\_defaut\\ score\_points\\ seuil\_decision\\ decision\\ shap\_top\_feature\\
    shap\_top\_valeur\\ date\_prediction\\[2pt] {\scriptsize UNIQUE (pret\_id, version\_modele)}};
  \node[table, anchor=north] (echa) at (0,-7.3) {\textcolor{white}{\textbf{echantillons\_modele}}\nodepart{two}
    \textbf{\textit{pret\_id}} \hfill PK, FK\\ echantillon};
  \node[table, anchor=north, draw=grisTexte, rectangle split part fill={grisTexte, white}] (par) at (0,-10.2)
    {\textcolor{white}{\textbf{parametres\_risque}}\nodepart{two}
    \textbf{cle} \hfill PK\\ valeur\\ valeur\_date\\ description};
  \node[tabled, anchor=north] (jou) at (11.0,-7.3) {\textcolor{white}{\textbf{journal\_statuts}} {\scriptsize\color{white}(ch.~1)}\nodepart{two}
    \textbf{journal\_id} \hfill PK\\ \textit{pret\_id} \hfill FK\\ ancien\_statut, nouveau\_statut\\
    jours\_retard, date\_effet\\ enregistre\_le, enregistre\_par};
  \node[tabled, anchor=north] (prov) at (11.0,-10.9) {\textcolor{white}{\textbf{provisions\_mensuelles}} {\scriptsize\color{white}(ch.~1)}\nodepart{two}
    \textbf{date\_arrete, \textit{pret\_id}} \hfill PK\\ jours\_retard, stage\_ifrs9\\ ead, pd\_utilisee\\ lgd, ecl};
  % liens
  \draw[->, thick, bleuML] ([yshift=-0.75cm]cli.north east) -- node[card, pos=0.15] {1} node[card, pos=0.85] {n} ([yshift=-0.75cm]pret.north west);
  \draw[<-, thick, bleuML] ([yshift=-0.75cm]ech.north west) -- node[card, pos=0.15] {n} node[card, pos=0.85] {1} ([yshift=-0.75cm]pret.north east);
  \draw[<-, thick, bleuML] (pred.north) -- node[card, pos=0.2] {n} node[card, pos=0.8] {1} (pret.south);
  \draw[<-, thick, bleuML] (echa.north east) -- node[card, pos=0.12] {0..1} node[card, pos=0.85] {1} (pret.south west);
  \draw[<-, thick, grisTexte, dashed] (jou.north west) -- node[card, pos=0.12] {n} node[card, pos=0.85] {1} (pret.south east);
  \draw[<-, thick, grisTexte, dashed] (prov.west) -- ++(-0.45,0) node[card, pos=0.5] {n} |- ($(pret.south east)+(-0.35,-0.5)$) -- ($(pret.south east)+(-0.35,0)$);
\end{tikzpicture}
\caption{Modèle relationnel de \code{risk_analytics_db}. En gras : clé primaire ; en italique : clé
étrangère. Chaque flèche part de la table référencée et pointe vers la table qui porte la clé étrangère
(le côté « n »). En pointillés : les deux tables que tu créeras dans
l'exercice 1.3, qui référencent toutes deux \code{demandes_pret}.}
\label{fig:mld}
\end{figure}

% -----------------------------------------------------------------------------
\section{Les choix de conception, et pourquoi}

Un recruteur ne te demandera pas de réciter le schéma. Il te demandera \textbf{pourquoi} tu l'as conçu
ainsi. Voici les huit décisions qui font de cette base une base « d'entreprise ».

\subsection*{1. NUMERIC pour l'argent, jamais FLOAT}
Un \code{FLOAT} stocke les nombres en binaire et ne représente pas exactement 0,10. Dans PostgreSQL :

\begin{lstlisting}[style=sql]
SELECT 0.1::float8 + 0.2::float8 = 0.3::float8 AS float_egal,     -- false
       0.1::numeric + 0.2::numeric = 0.3       AS numeric_egal,   -- true
       (0.1::float8 + 0.2::float8)::text       AS float_valeur;   -- 0.30000000000000004
\end{lstlisting}
Sur un million de mensualités, ces erreurs s'accumulent et les totaux comptables ne tombent plus juste.
\code{NUMERIC(12,2)} calcule en décimal exact, au centime.

\subsection*{2. Le retard n'est pas stocké}
Le DPD change chaque jour sans événement (partie 2). Il est calculé dans une vue à partir des dates.
Même raisonnement pour l'encours : on le déduit des paiements.

\subsection*{3. ON DELETE RESTRICT sur toutes les clés étrangères}
Supprimer un client ne doit \textbf{jamais} effacer silencieusement ses prêts et leur historique : les
banques ont des obligations de conservation. Avec \code{RESTRICT}, la suppression est refusée tant que
des prêts y font référence. \code{CASCADE} serait dangereux ici.

\subsection*{4. Des identifiants générés par la base}
\code{GENERATED BY DEFAULT AS IDENTITY} pour les clients et les prêts : la base attribue un numéro,
mais on peut fournir le sien au chargement (on garde le numéro de ligne du fichier source, ce qui relie
chaque prêt à tes sorties R). \code{GENERATED ALWAYS} pour les échéances et les prédictions : personne
ne doit choisir ces numéros à la main.

\subsection*{5. Une table de paramètres, source unique de vérité}
La LGD (45~\%), les seuils de 30 et 90 jours et la date de référence sont stockés dans
\code{parametres_risque}, pas écrits en dur dans dix requêtes. Si le seuil change, on modifie une ligne.
Une contrainte impose qu'un paramètre ait soit une valeur numérique, soit une date, jamais les deux.

\subsection*{6. Des contraintes entre colonnes}
Certaines règles portent sur plusieurs colonnes d'une même ligne. Par exemple, une décision doit être
cohérente avec la PD et le seuil : \code{(decision = 'refuser') = (probabilite_defaut >= seuil_decision)}.
L'égalité entre deux conditions exprime une équivalence (« si et seulement si »). On verra cette astuce
en détail au chapitre~1.

\subsection*{7. Un registre versionné des prédictions}
Plusieurs prédictions par prêt (une par version de modèle), jamais de mise à jour, un drapeau
\code{est_champion} pour le modèle qui décide, le seuil enregistré avec la prédiction.

\subsection*{8. Deux statuts distincts}
\code{statut_dossier} décrit la \textbf{décision} (en attente, accordé, refusé) ;
\code{statut_pret} décrit le \textbf{remboursement} (sain, en retard, en défaut, soldé). Le second n'a de
sens que si le premier vaut « accordé » : une contrainte l'impose.

% -----------------------------------------------------------------------------
\section{Dictionnaire de données}

Pour chaque table : les colonnes, leur type PostgreSQL, leurs contraintes et leur signification.

\newcommand{\dico}[1]{%
  \par\needspace{6\baselineskip}\noindent\textbf{\large\code{#1}}\par\vspace{2pt}}
\newenvironment{tabdico}{\noindent\footnotesize\tabularx{\linewidth}{@{}P{3.35cm} P{2.4cm} P{4.05cm} L@{}}
  \toprule \textbf{Colonne} & \textbf{Type} & \textbf{Contraintes} & \textbf{Description} \\ \midrule}
  {\bottomrule\endtabularx\par\medskip}

\dico{clients} {\small Le référentiel des emprunteurs. 32\,409 lignes (+ Awa Koné).}\par
\begin{tabdico}
\code{client_id} & BIGINT & PK, identité & identifiant de l'emprunteur \\
\code{nom_complet} & VARCHAR(100) & NOT NULL & nom fictif (Faker) \\
\code{email} & VARCHAR(120) & NOT NULL, UNIQUE & adresse fictive \code{@exemple.fr} \\
\code{age} & SMALLINT & NOT NULL, entre 18 et 100 & âge à la demande \\
\code{revenu_annuel} & NUMERIC(12,2) & NOT NULL, $> 0$ & revenu annuel en euros \\
\code{statut_logement} & VARCHAR(20) & NOT NULL, liste fermée & locataire, proprietaire, credit\_immobilier, autre \\
\code{anciennete_emploi} & NUMERIC(4,1) & $\geq 0$, $\leq$ âge $-$ 14 & années dans l'emploi actuel ; NULL si inconnue (887 clients) \\
\code{anciennete_credit} & SMALLINT & NOT NULL, $\geq 0$ & années d'historique de crédit \\
\code{defaut_anterieur} & BOOLEAN & NOT NULL & un défaut a déjà été enregistré par le passé \\
\code{date_inscription} & DATE & NOT NULL & entrée dans la banque \\
\end{tabdico}

\dico{demandes_pret} {\small Le dossier de crédit, la décision et le statut de remboursement.}\par
\begin{tabdico}
\code{pret_id} & BIGINT & PK, identité & numéro du prêt (= ligne du fichier source) \\
\code{client_id} & BIGINT & NOT NULL, FK $\rightarrow$ clients, RESTRICT & l'emprunteur \\
\code{montant_demande} & NUMERIC(12,2) & NOT NULL, $> 0$ & capital emprunté \\
\code{duree_mois} & SMALLINT & NOT NULL, entre 6 et 120 & durée du prêt (12 à 60 mois dans les données) \\
\code{taux_interet} & NUMERIC(5,2) & NOT NULL, $0 < t < 40$ & taux annuel en \% \\
\code{taux_impute} & BOOLEAN & NOT NULL, défaut false & taux absent de la source, remplacé par la médiane (3\,094 prêts) \\
\code{motif_pret} & VARCHAR(25) & NOT NULL, liste fermée & etudes, sante, creation\_entreprise, personnel, travaux, regroupement\_dettes \\
\code{note_risque} & CHAR(1) & NOT NULL, de A à G & note interne (A = meilleur risque) \\
\code{date_demande} & DATE & NOT NULL & dépôt du dossier \\
\code{date_decision} & DATE & $\geq$ \code{date_demande} & décision d'octroi \\
\code{date_premiere_echeance} & DATE & & première mensualité \\
\code{statut_dossier} & VARCHAR(12) & NOT NULL, défaut en\_attente & en\_attente, accorde, refuse \\
\code{statut_pret} & VARCHAR(12) & non NULL $\Leftrightarrow$ accordé & sain, en\_retard, en\_defaut, solde \\
\end{tabdico}

\dico{echeancier_remboursements} {\small Une ligne par mensualité, avec le suivi du paiement.
1\,107\,936 lignes.}\par
\begin{tabdico}
\code{echeance_id} & BIGINT & PK, identité ALWAYS & identifiant technique \\
\code{pret_id} & BIGINT & NOT NULL, FK $\rightarrow$ demandes\_pret & le prêt \\
\code{num_echeance} & SMALLINT & NOT NULL, $> 0$ ; UNIQUE avec \code{pret_id} & rang de la mensualité (1 à $n$) \\
\code{date_echeance} & DATE & NOT NULL & date à laquelle la mensualité est due \\
\code{montant_du} & NUMERIC(12,2) & NOT NULL, $> 0$ & mensualité = capital + intérêts \\
\code{part_capital} & NUMERIC(12,2) & NOT NULL, $\geq 0$ & part qui rembourse la dette \\
\code{part_interets} & NUMERIC(12,2) & NOT NULL, $\geq 0$ & rémunération de la banque \\
\code{capital_restant_du} & NUMERIC(12,2) & NOT NULL, $\geq 0$ & dette restante \emph{après} cette échéance \\
\code{montant_paye} & NUMERIC(12,2) & NOT NULL, défaut 0, $\leq$ \code{montant_du} & somme déjà versée \\
\code{date_paiement} & DATE & non NULL $\Leftrightarrow$ entièrement payée & date du paiement complet \\
\end{tabdico}

\dico{predictions_scoring_ml} {\small Le registre des prédictions : 3 modèles $\times$ 32\,409 prêts.}\par
\begin{tabdico}
\code{prediction_id} & BIGINT & PK, identité ALWAYS & identifiant technique \\
\code{pret_id} & BIGINT & NOT NULL, FK & le prêt évalué \\
\code{version_modele} & VARCHAR(40) & NOT NULL ; UNIQUE avec \code{pret_id} & grille\_score\_v1, logistique\_v1, xgboost\_monotone\_v2 \\
\code{est_champion} & BOOLEAN & NOT NULL, défaut false & vrai pour le modèle qui décide en production \\
\code{probabilite_defaut} & NUMERIC(7,6) & NOT NULL, entre 0 et 1 & PD prédite \\
\code{score_points} & SMALLINT & & points de la grille de score (grille uniquement) \\
\code{seuil_decision} & NUMERIC(4,3) & NOT NULL, entre 0 et 1 & seuil de refus du modèle \\
\code{decision} & VARCHAR(10) & NOT NULL ; refuser $\Leftrightarrow$ PD $\geq$ seuil & accorder ou refuser \\
\code{shap_top_feature} & VARCHAR(50) & & variable SHAP la plus influente (champion uniquement) \\
\code{shap_top_valeur} & NUMERIC(8,4) & & sa contribution (log-odds) \\
\code{date_prediction} & TIMESTAMP & NOT NULL & date et heure du calcul \\
\end{tabdico}

\dico{echantillons_modele} {\small Le découpage du projet R : 25\,927 prêts d'apprentissage, 6\,482 de test.}\par
\begin{tabdico}
\code{pret_id} & BIGINT & PK, FK $\rightarrow$ demandes\_pret & le prêt \\
\code{echantillon} & VARCHAR(13) & NOT NULL, apprentissage ou test & échantillon d'appartenance \\
\end{tabdico}

\dico{parametres_risque} {\small Les paramètres métier (4 lignes).}\par
\begin{tabdico}
\code{cle} & VARCHAR(40) & PK & lgd, seuil\_defaut\_jours, seuil\_stage2\_jours, date\_reference \\
\code{valeur} & NUMERIC & une seule des deux & valeur numérique (0,45 ; 90 ; 30) \\
\code{valeur_date} & DATE & une seule des deux & valeur de type date (30/09/2026) \\
\code{description} & TEXT & NOT NULL & signification et source réglementaire \\
\end{tabdico}

Les tables \code{journal_statuts} et \code{provisions_mensuelles} sont spécifiées dans l'exercice 1.3 :
tu les concevras toi-même.

% -----------------------------------------------------------------------------
\section{Les index B-Tree}

Un index est une structure annexe qui permet de retrouver des lignes sans lire toute la table (on
verra son fonctionnement au chapitre~1). PostgreSQL crée automatiquement un index pour chaque clé
primaire et chaque contrainte \code{UNIQUE}, mais \textbf{pas pour les clés étrangères}.

\begin{center}\footnotesize
\begin{tabularx}{\linewidth}{@{}P{4.4cm} P{3.9cm} L@{}}
\toprule
\textbf{Index} & \textbf{Colonnes} & \textbf{À quoi il sert} \\
\midrule
\multicolumn{3}{@{}l}{\emph{Créés automatiquement}} \\
clés primaires (\code{*_pkey}) & l'identifiant de chaque table & retrouver une ligne par son identifiant \\
\code{clients_email_key} & \code{email} & garantir l'unicité de l'adresse \\
\code{une_echeance_par_numero} & \code{(pret_id, num_echeance)} & unicité, et toutes les échéances d'un prêt (33~Mo) \\
\code{une_prediction_par_version} & \code{(pret_id, version_modele)} & unicité, et les prédictions d'un prêt \\
\midrule
\multicolumn{3}{@{}l}{\emph{Créés par le script}} \\
\code{idx_prets_client} & \code{client_id} & jointure client $\rightarrow$ prêts (clé étrangère) \\
\code{idx_prets_statut} & \code{statut_pret} & filtrer un statut rare (933 prêts en retard) ; pour un statut
fréquent, lire la table est plus rapide \\
\code{idx_echeances_impayees} & \code{(pret_id, date_echeance)} \textbf{partiel} : \code{WHERE montant_paye < montant_du} &
trouver l'impayé le plus ancien d'un prêt : le calcul du retard. 13~Mo seulement, car il ne contient
que les 442\,459 échéances non soldées \\
\bottomrule
\end{tabularx}
\end{center}

% -----------------------------------------------------------------------------
\section{Le code DDL complet}

Le script ci-dessous crée toute la base. Il est fourni dans \code{sql/01_schema.sql}. Lis-le en entier
en t'appuyant sur le dictionnaire : chaque commentaire justifie un choix. Il commence par supprimer les
tables existantes, ce qui permet de le relancer pour repartir de zéro.

\lstinputlisting[style=long]{../sql/01_schema.sql}

\begin{resultat}[Comment l'exécuter et ce que tu dois obtenir]
Dans un terminal, depuis le dossier du projet :
\begin{lstlisting}[style=shell]
psql -d risk_analytics_db -v ON_ERROR_STOP=1 -f sql/01_schema.sql
\end{lstlisting}
L'option \code{ON_ERROR_STOP=1} arrête le script à la première erreur au lieu de continuer. Dans
DBeaver : ouvre le fichier, puis \emph{Exécuter le script} (\code{Alt+X}). La commande \code{\dt} de
psql doit lister 6 tables : \code{clients}, \code{demandes_pret}, \code{echantillons_modele},
\code{echeancier_remboursements}, \code{parametres_risque}, \code{predictions_scoring_ml}.
\end{resultat}

\begin{resume}
\begin{enumerate}
\item On conçoit en trois temps : modèle conceptuel (entités, associations, cardinalités), modèle
relationnel (tables, clés), code DDL.
\item Une association « un à plusieurs » devient une clé étrangère du côté « plusieurs ».
\item Les choix clés : \code{NUMERIC} pour l'argent, retard calculé et non stocké, \code{ON DELETE
RESTRICT}, identités, table de paramètres, contraintes entre colonnes, registre versionné, deux statuts.
\item Les clés primaires et \code{UNIQUE} sont indexées automatiquement, pas les clés étrangères. Un
index partiel n'indexe que les lignes utiles.
\end{enumerate}
\end{resume}
